\section{Introduction}

The contest asks the participant to schedule a memory-bounded DAG of
MatMul and Pointwise operators on a two-level memory hierarchy, using
only the Gemini API for reasoning. A schedule makes four interacting
decisions per subgraph: \emph{subgraphing}, \emph{retention},
\emph{granularity} $[w,h,k]$, and \emph{traversal order}. The
structural observation driving our design is that partitioning
dominates the complexity, while granularity and traversal can be
searched locally with small enumerations.

Our system splits the two classes of reasoning between the LLM and
deterministic code. Gemini handles the combinatorial, semantically
rich task of proposing a partition and retention plan; Python code
handles granularity and traversal search, where exhaustive
enumeration over a small candidate set is both feasible and precise.
The agent runs a best-of-$N$ loop over Gemini attempts with a
cost-feedback prompt, and closes the loop with a long final
grain-refinement pass over the best partition found.


\section{Challenges for an LLM Partition Designer}

Three properties of the partition space are difficult for a language model.
(1)~\emph{Irregular cost}: small partition edits cause large discontinuous
jumps in memory traffic or feasibility, so the LLM cannot reliably improve
a partition by incremental rewrites. (2)~\emph{Hidden arithmetic}: whether
an intermediate becomes ephemeral depends on consumer sets across the
whole subgraph; whether a given $[w,h,k]$ fits fast memory depends on the
per-role working-set formula (LHS resident $h\times K$ for chain-upstream
MM, etc.). These are tedious and error-prone for a text-generating model.
(3)~\emph{Stochastic sampling}: the same prompt yields different partitions
across runs.

We address (1) and (2) by pre-digesting the graph for the LLM and by
letting deterministic code enforce every hard constraint; (3) with a
best-of-$N$ harness and a cost-feedback log threaded through the
prompts.


% ════════════════════════ §3 System ════════════════════════════════
\section{System Overview}

Each invocation proceeds through four stages.

\paragraph{1.~Safety-net fallback.}
Emit a 1-op-per-subgraph solution with an analytic grain picker and
write it to disk before any Gemini call. Guarantees a valid output if
later stages fail.

\paragraph{2.~Fallback refinement.}
Run 2-stage grain search over the fallback partition using $\sim$30\%
of the time budget. This alone closes 5--15\% of the gap on most
benchmarks.

\paragraph{3.~Best-of-$N$ Gemini loop.}
Three attempts with distinct (strategy, temperature) pairs:
\texttt{(rich, 0.25)}, \texttt{(minimal, 0.70)}, \texttt{(rich, 0.55)}.
Each attempt proposes a partition and retention plan; our code runs
grain search over it and keeps the best valid solution. A compact
cost-feedback summary of prior attempts is threaded into the next
prompt so Gemini can reason about what structural change to try next.

\paragraph{4.~Final refinement.}
After the Gemini loop, spend all remaining budget on a high-quality
grain search over the best partition seen. During Gemini attempts the
per-subgraph grain budget is necessarily short (2--3~s for the 100-op
benchmark), often too short to reach the optimum; the final pass gives
each subgraph 10--15~s of dedicated search.


% ════════════════════════ §4 Prompts ════════════════════════════════
\section{Prompt Design}

The system prompt encodes every evaluator rejection path as six
invariants: coverage and producer-before-consumer order; retention
validity (only op-produced tensors can be retained, lifetime one
step); the ephemeral-with-external-consumer rule; the unified
execution grid, including the clarification that ephemeral
intermediates are exempt from the shape-matching requirement; the
MM-chain fusion rules under full-K vs.\ split-K; and the per-role
working-set formula (LHS-head resident $h\times K$, standalone MM
LHS $h\times k_\mathrm{eff}$, pointwise $w\times h$, etc.). Encoding
these invariants in the prompt means Gemini knows, before proposing
anything, exactly which configurations will be rejected.

The user prompt carries a client-side pre-analysis of the graph so
Gemini does not have to scan the raw JSON to find structure:
output-shape buckets, shared graph-input clusters, pre-checked
fusion-candidate edges, and \emph{triple-fusion candidates} --- ordered
triples of the form (pointwise, matmul, pointwise) where the matmul's
two adjacent intermediates can both be made ephemeral. Each triple is
annotated with a producer-shape \emph{masking warning} when an earlier
MM producing the triple's input has a larger output shape, to prevent
Gemini from extending the fusion backwards into a large-shape MM
--- which would otherwise force the subgraph to tile at the small
output's grid and leave the large MM mostly masked.


% ════════════════════════ §5 Triple fusion ═════════════════════════
\section{Triple Fusion with Ephemeral MatMul Operands}

The single biggest improvement came from Gemini learning to fuse a
three-operator chain: an upstream pointwise whose output feeds the
MatMul's LHS, the MatMul itself, and a downstream pointwise that
consumes the MatMul's output. Because both of the MatMul's adjacent
tensors are produced \emph{and} consumed entirely inside the subgraph
(and neither is listed in \texttt{tensors\_to\_retain}), both become
ephemeral by definition. Crucially, the usual dominant working-set
term --- the MatMul's resident LHS of size $h \times K$ --- \textbf{%
disappears} because the LHS is no longer a resident input: it is
produced tile-by-tile by the upstream pointwise and consumed
immediately by the MatMul.

For a large-$K$ operator (e.g.\ $K = 4096$ with the per-benchmark
capacity of a few hundred thousand cells), this shrinks the working
set enough to permit full-$K$ execution ($nk = 1$), eliminating the
split-$K$ reload penalty that dominates the per-op baseline. Once the
system prompt includes this working-set calculation explicitly,
Gemini identifies the triples in the graph and proposes the fusion
directly.


% ════════════════════════ §6 Grain search ═══════════════════════════
\section{Grain Search}

Candidates are a union of powers-of-two, divisors of the output
dimension, the native granule, and a \emph{tile-count sweep} so
non-standard values are reachable (e.g.\ $w = 103$ for
$W_\mathrm{out} = 4096$ with 40 tiles, where neither $103$ nor its
complement is a power of two). Snake traversals (row-snake and
column-snake) are only useful when the reused operand matches the
unchanging tile axis, and the benefit is multiplicative on top of the
row-major cost, so we use a two-stage search: Stage~1 enumerates all
candidates under row-major traversal and keeps the top-10 by cost;
Stage~2 tries the two snake variants on that shortlist. This is
roughly $3\times$ faster than evaluating all three traversals at
every candidate --- necessary to keep the 103-op benchmark inside the
time budget.

The tile grid on which each snake permutation is built is derived
from the \emph{non-ephemeral} outputs of the subgraph, matching the
evaluator's own grid definition. A subgraph containing an ephemeral
4096$\times$1024 intermediate and a non-ephemeral 1024$\times$1024
output tiles over a $1024\times 1024$ grid, not the $4096\times 1024$
one, so the snake permutations must have length $\lceil 1024/w
\rceil \cdot \lceil 1024/h \rceil$.


% ════════════════════════ §8 Reproducibility ═══════════════════════
\section{Reproducibility}

Submission contains \texttt{agent.py} (entry point),
\texttt{prompts/} (system and repair templates), and
\texttt{requirements.txt} (\texttt{google-genai}). Invocation is the
contest-specified
\texttt{python3 agent.py <problem.json> <output.json>}; the agent
expects \texttt{GOOGLE\_API\_KEY} in the environment and opens no
other network connection. The \texttt{TEAMKOREA\_BUDGET\_S} environment
variable (default 540~s) leaves a 60-second margin inside the
organiser's 600-second per-benchmark timeout.
